
\section{Analysis on Shard Synchronization Network Overhead}

Define
\begin{itemize}
    \item $n$ as the number of nodes
    \item $c$ as the number of shard copies preserved by the active layer
    \item $d$ as the number of shards
    \item $d_t$ as the average number of shards accessed by each transaction
    % \item $S$ as the state size, $s=\frac{S}{d}$ as the shard size
    \item $t$ as the transaction and reply size
    \item $\sigma$ as the missing rate of local shard cache
\end{itemize}

The expectation of the network overhead of synchronizing shards for each transaction would be
\[
    d_t(n-c)s\sigma
\]
and the network overhead of dispersing transactions and replying to clients would be
\[
    nt
\]
note that although BFT protocols usually requires only $f+1$ replies to client, all nodes are sending replies, so the network overhead is still incurred.

Compared to the baseline overhead for dispersing transactions and replying to clients, shard synchronization makes a
\[
    1+\frac{n-c}{n}\cdot{}d_t\cdot{}\frac{s}{t}\sigma
\]
times amplification of network overhead.

Let's break down analysis each of the fractions in the context of a simple key-value store state
\begin{itemize}
    \item $\frac{n-c}{n}$ increases as $n$ increases (because $c$ is constant) but never reach $1$.
    \item $d_t$ is the average number of shards access by the transactions, expected to be only slightly greater than $1$.
    \item $\frac{s}{t}$ has varying meanings with different states and transactions.
    With key-value store state it is approximately the number of key-value pairs in a shard.
    \item $\sigma$ is the missing rate which is also transaction-dependent.
    Given a skewed workload (\eg most configurations in YCSB~\cite{cooper2010benchmarking}), the missing rate can be as low as 1-10\%.
\end{itemize}

We perform the numeric analysis with $d_t=1.5$, $\frac{s}{t}=10$ and $\sigma=0.1$ and the result is shown in the paper.
\xxx{Derive these parameters from real world workloads.}